\subsection{ISIC Prompt Ablation}

Using the repository's Med-SA adapter implementation, we trained four prompt
configurations on the available ISIC 2016 dataset: one positive click, three
positive clicks, BBox 0.5, and BBox 0.75. All runs used a SAM ViT-B checkpoint,
1024-pixel inputs, batch size 4, and 40 epochs. Metrics below are from final
validation at epoch 39.

\begin{table*}[t]
    \centering
    \caption{Melanoma segmentation results: replication on ISIC 2016 compared with the Med-SA paper results on ISIC 2019.}
    \label{tab:melanoma_prompt_comparison}
    \begin{tabular}{lcccc}
        \hline
        & \multicolumn{2}{c}{Replication (ISIC 2016)}
        & \multicolumn{2}{c}{Med-SA paper (ISIC 2019)} \\
        Prompt & Dice & IoU & Dice & IoU \\
        \hline
        1 point     & 91.7 & 85.6  & 92.6 & 84.1 \\
        3 points    & 91.7 & 85.53 & 93.4 & 84.7 \\
        BBox 0.5    & 94.8 & 90.41 & 92.1 & 83.0 \\
        BBox 0.75   & 95.4 & 91.4  & 93.0 & 84.2 \\
        \hline
    \end{tabular}
\end{table*}

BBox 0.75 performed best in this experiment, exceeding three-click prompting
by 5.87 IoU points and 3.7 Dice points based on the rounded values in
Table~\ref{tab:melanoma_prompt_comparison}. The BBox 0.5 result was also higher
than the three-click result. These findings suggest that, on this dataset and
under this prompt-generation procedure, box prompts were more effective than
three positive clicks. They do not establish that boxes are generally superior:
each configuration was run once, and prompt construction can materially affect
the comparison.

\paragraph{Prompt protocol.}
The paper reports four prompt conditions: one positive point, three positive
points, and bounding boxes with 50\% or 75\% overlap with the target. The
repository's ISIC path originally generated one random foreground click and
did not expose these alternatives as command-line settings. The prompt
ablation therefore required extending prompt generation and routing the
selected point or box to SAM. In particular, ``BBox 0.5'' and ``BBox 0.75''
are prompt-generation conditions, not instructions to multiply box coordinates
by 0.5 or 0.75. Our generated boxes were derived from the ground-truth masks;
this simulates an annotation-provided prompt and is not an inference of the
box by the model. Since the paper does not give a fully specified box
construction algorithm in the released implementation available here, our
box-generation procedure may not match the authors' exactly. Likewise, our
three clicks were sampled randomly as positive foreground points; they should
not be described as reproducing an iterative correction procedure unless each
later click was selected from the model's preceding segmentation errors.

\paragraph{Dataset and mask availability.}
The comparison is not an exact reproduction because our runs used ISIC 2016,
whereas the paper reports melanoma results on ISIC 2019, and we used ViT-B
rather than the paper's ViT-H setup. The image collection being public does
not, by itself, provide the pixel-level targets needed to train and score a
segmentation model. The official ISIC catalogue's 2016 dataset table includes a Task 3B entry with
900 training images and associated binary masks, as well as 379 test images
and associated masks~\cite{isicdata}. These are masks for the 2016 release;
they are not evidence that the pixel-level annotations used for the paper's
ISIC 2019 experiment are available in the same bundle. The paper and the
released repository do not identify a specific 2019 segmentation-mask
release, provide the corresponding image-ID manifest, or specify the exact
train/test split used to generate the reported melanoma metrics. Therefore,
the exact ISIC 2019 segmentation results cannot be independently reproduced
from the paper's dataset name and released code alone: without the matching
ground-truth masks and split, neither supervised training nor segmentation
scoring can be reconstructed. This is an artifact and protocol limitation,
not a claim that no ISIC masks exist for any year. The repository's checked-in
ISIC loader expects ISBI 2016 Part 1 filenames, so it supports the available
2016 mask bundle but requires an explicit loader or data-layout change for a
different release.

\paragraph{Interpretation.}
The results support a within-experiment comparison of prompt modes on the
available ISIC 2016 data. They should not be treated as directly comparable
to the paper's ISIC 2019 numbers because the dataset and backbone differ, the
box-generation procedure may differ, and the paper's mask release and split
are not identified in the released artifacts. The study also includes no repeated seeds or
uncertainty estimates. Accordingly, this is a limited replication: it shows
that the tested box prompts performed well in this implementation, but it
does not establish reproduction of the published ISIC 2019 result.

\begin{thebibliography}{1}
\bibitem{isicdata}
International Skin Imaging Collaboration,
\emph{ISIC Challenge Datasets},
\url{https://challenge.isic-archive.com/data/}.
Accessed October 8, 2026.
\end{thebibliography}